\chapter{Methodology}\label{Ch:Methodology}
%=====================================================================================%
% Chapter outline:  
%       - Flumy
%           - Modelling rules
%           - Parameter ranges
%       - GAN framework
%           - Model architecture(s)
%           - Conditioning approach
%           - (Hyperparameters)
%       - Evaluation framework
%           - MSSWD
%           - 
%           - Well data
%           - (Heat)flow simulations
%=====================================================================================%
This chapter aim to clarify the complete workflow of this project. This includes the generation on 3D datasets used for training, validating and testing the model which is covered in section \ref{Sec:data_gen}. Section \ref{Sec:GAN_framew} described the exact GAN framework used in the project. This includes the facies representation approach as well as the validation function and loss function which are used to monitor the model during training. Finally, section \ref{Sec:eval} introduces the evaluation metrics used to test geological realism in the final realizations.
\section{Training data generation}\label{Sec:data_gen}
In order to train the Deep learning algorythm a large number of training samples representing the ground truth ar needed. Since the subsurface is a 3D domain with a high degree of uncertainty the current method to obtaining a large number of training samples is by simply generating them using a quick algorythm (\cite{rongier2025atowards}; \cite{sun2023geological}). The approach taken in this project is to start of with three different geological conceptual models, two of which are based on the upper and lower plain delta environment and the final conceptual model has a general deltaic representation. The 

\section{GAN framework}\label{Sec:GAN_framew}

\section{Conditioning to well data}\label{Sec:Inversion}

\section{Evaluation Framework}\label{Sec:eval}